\begin{textsample}{NEG Dim 1 – ap – Score 3.00 – ap/ap\_ef\_25.txt}  \label{ex:f1_neg_117}
4.4.42.1.
Educação Física nos Anos Finais do Ensino Fundamental : Unidades Temáticas, Objetos de Conhecimento e Habilidades De acordo com a BNCC no Ensino Fundamental – Anos Finais, é de se esperar que os estudantes estejam preparados para se deparar com diversos docentes, que atuem com maior independência e dominem uma série de conhecimentos, o que torna mais complexas as interações e a sistemática de estudos.
Ainda assim, os alunos nessa fase de escolarização têm maior capacidade de apreciação das mais diversas manifestações da cultura corporal, podendo ocorrer com a incorporação de mais aspectos e detalhes.
Ao vivenciá -las, os alunos podem apreciar a beleza, a estética, discutir o contexto de sua produção, avaliar algumas técnicas e estratégias, observar os padrões de movimento, entre inúmeras outras possibilidades.
Podem, principalmente, aprender a contemplar essa diversidade e perceber as inúmeras opções que existem, tanto para praticar quanto para apreciar ( PCNs, 1998 ).
Essas características permitem aos estudantes maior aprofundamento nos estudos das práticas corporais na escola.
Nesse contexto, e para aumentar a flexibilidade na delimitação dos currículos e propostas curriculares, tendo em vista a adequação às realidades locais, as habilidades de Educação Física para o Ensino Fundamental – Anos Finais foram organizadas em dois blocos ( 6º e 7º anos ; 8º e 9º anos ) e se referem às unidades temáticas e seus respectivos objetos de conhecimento.

% matched lemmas: 
\end{textsample}
